\chapter{conjuntos, relaciones y funciones}


\textbf{Relaciones:}

\begin{definicion}[title=Definicion]
    \textbf {si A y B son dos conjuntos, una relación $R$ de $A$ en $B$ es un subconjunto cualquiera R del producto cartesiano $A \times B$. Es decir, $R \subseteq A \times B$ o $R \in P (A \times B)$.}
\end{definicion}

\begin{teoria}[title=Propiedades de las relaciones]
\begin{listapropiedades}
  \item[Reflexiva] $R$ es reflexiva si $\forall a \in A,\ (a,a)\in R$. Dicho de otra manera , R es reflexiva si en cada vertice hay una flecha "bucle".
  \item[Simétrica] $R$ es simétrica si $\forall a,b \in A,\ (a,b)\in R \Rightarrow (b,a)\in R$. Es decir, $R$ es simétrica si para cada flecha que va de $a$ a $b$, hay una flecha que va de $b$ a $a$.
  \item[Antisimétrica] $R$ es antisimétrica si $\forall a,b \in A,\ (a,b)\in R \wedge (b,a)\in R \Rightarrow a=b$.
  \item[Transitiva] $R$ es transitiva si $\forall a,b,c \in A,\ (a,b)\in R \wedge (b,c)\in R \Rightarrow (a,c)\in R$.
\end{listapropiedades}
\end{teoria}


\begin{teoria}[title=Relación de equivalencia y relación de orden]
  Sea A un conjunto y $R$ una relación en A. 
  \begin{listapropiedades}
    \item[Relación de equivalencia] sea A un conjunto y $R$ una relación de A en A. 
      Se dice que R es una relación de equivalencia si R es reflexiva, simétrica y transitiva.
    \item[Relación de orden] sea A un conjunto y $R$ una relación de A en A. 
      Se dice que R es una relación de orden si R es reflexiva, antisimétrica y transitiva.
  \end{listapropiedades}
\end{teoria}

%\begin{proposicion}
%  Para todo $a, b \in \mathbb{R}$, se cumple que $(a+b)^2 = a^2 + 2ab + b^2$.
%\end{proposicion}

%\begin{demostracion}
%  Desarrollando el producto $(a+b)(a+b)$ y aplicando la propiedad distributiva...
%\end{demostracion}